\runninghead{The Sunday and Its Texts}
\section{The Last Sunday after Pentecost and Its Texts}
\label{sec:opening}

The Gospel of this Mass is taken from the discourse that Jesus gave his
disciples on the Mount of Olivet. Leaving the temple, he had told them that
not a stone of its buildings would be left upon a stone; when he sat down on
the mountain they came to him privately and asked, ``Tell us when shall these
things be? And what shall be the sign of thy coming and of the consummation of
the world?'' (Mt 24:1--3). The Mass begins its reading at the fifteenth verse
of his answer. When they see ``the abomination of desolation, which was spoken
of by Daniel the prophet, standing in the holy place,'' those in Judea are to
flee to the mountains; the man on the housetop is not to come down for
anything in his house, nor the man in the field to go back for his coat; there
is woe to women with child and nursing, and they are to pray that their flight
be not in winter or on the sabbath. The tribulation will be such as has never
been, and the days will be shortened for the sake of the elect. False christs
and false prophets will show great signs; if anyone says that Christ is in the
desert or in the inner rooms, they are not to believe it, for the coming of the
Son of man will be like lightning from the east to the west, and ``Wheresoever
the body shall be, there shall the eagles also be gathered together.'' After
the tribulation the sun and moon will be darkened, the stars will fall, the
sign of the Son of man will appear in heaven, all the tribes of the earth will
mourn and see him ``coming in the clouds of heaven with much power and
majesty,'' and his angels will gather his elect from the four winds with a
trumpet. The reading ends with a parable from the fig tree, whose tender
branch and new leaves show that summer is near, and with two sayings: ``this
generation shall not pass till all these things be done,'' and ``Heaven and
earth shall pass: but my words shall not pass'' (Mt 24:15--35).

The Epistle is the prayer with which St Paul opens his letter to the
Colossians, a church he knew through the report of Epaphras, who had taught
them (Col 1:4, 7--8). Paul and Timothy have not ceased to pray that the
Colossians may be filled with the knowledge of the will of God ``in all wisdom
and spiritual understanding,'' that they may walk worthy of God, be fruitful in
every good work, and be strengthened ``in all patience and longsuffering with
joy,'' giving thanks to the Father ``who hath made us worthy to be partakers of
the lot of the saints in light: Who hath delivered us from the power of
darkness and hath translated us into the kingdom of the Son of his love, In
whom we have redemption through his blood, the remission of sins'' (Col
1:9--14).

The five chants are those of the Twenty-third Sunday after Pentecost, printed
out again in full. The Introit gives God's words in the letter Jeremias sent
from Jerusalem to the exiles carried away to Babylon (Jer 29:1, 4): ``The Lord
saith, I think thoughts of peace, and not of affliction; you shall call upon Me
and I will hear you; and I will bring back your captivity from all places,''
with the psalm verse ``Lord, thou hast blessed thy land: thou hast turned away
the captivity of Jacob'' (Jer 29:11--12, 14; Ps 84:2). In the letter those
words follow a warning against the prophets and diviners among the exiles and
the promise that God will visit them ``When the seventy years shall begin to
be accomplished in Babylon'' (Jer 29:8--10). The Gradual gives thanks:
``Thou hast delivered us, O Lord, from them that afflict us: and hast put them
to shame that hate us. In God we will glory all the day: and in Thy name we
will give praise for ever'' (Ps 43:8--9). The Alleluia and then the Offertory
sing ``From the depths I have cried to Thee, O Lord: Lord, hear my prayer''
(Ps 129:1--2), and the Offertory repeats its first line. The Communion is the
Lord's promise in Mark, ``whatsoever you ask when you pray, believe that you
shall receive, and it shall be done unto you'' (Mk 11:24), spoken when Peter
pointed out the fig tree that Jesus had cursed and that had withered from the
roots overnight. Jesus answered him, ``Have the faith of God,'' spoke of
faith that moves a mountain into the sea, gave this promise, and added at once
that those who stand to pray must forgive (Mk 11:20--25).

The three prayers of the priest belong to this Mass alone. The Collect asks
God to ``arouse \dots\ the wills of Thy faithful'' so that they ``may the more
abundantly partake of Thy mercies.'' The Secret asks him to ``accept the
offerings and prayers of Thy people'' and to ``turn all our hearts unto Thee,
that, being delivered from earthly desires, we may pass on to the enjoyments
of heaven.'' The Postcommunion asks that, ``through this sacrament which we
have received, whatever is evil in our hearts may be restored by its gift of
healing.''

The Missal prints this Mass at the end of the Proper of Time under the
heading \latin{Dominica vigesima quarta et ultima post Pentecosten}, ``the
twenty-fourth and last Sunday after Pentecost,'' of the second class (printed
pp.~422--424, nos.~1769--1778). It is said on the Sunday before the First
Sunday of Advent however many Sundays the year has. A note closing the
Twenty-third Sunday directs that, when there are more than twenty-four, the
Sundays left over after the Epiphany are taken up after the twenty-third,
\latin{Et ultimo loco semper ponitur Missa dominicae XXIV}, ``and in the last
place is always set the Mass of the twenty-fourth Sunday'' (printed p.~417).
In 2026 Easter fell on 5~April and Pentecost on 24~May, and there are
twenty-six Sundays after Pentecost. Under RG~18~b the Fifth and Sixth Sundays
after the Epiphany were said on 8 and 15~November, and this Mass is said on
22~November, the twenty-sixth Sunday by count, a week before the First Sunday
of Advent on 29~November. The calendar's feast of St Cecilia on 22~November is
of the third class and is not commemorated, because a Sunday of the second
class admits only the commemoration of a feast of the second class
(RG~111~b). The Mass is said in green (RG~127~b), with the Gloria and the
Credo, and with the Preface of the Most Holy Trinity, which the Missal directs
after the Secret (RGMR~494~b). It has a single Collect, Secret and
Postcommunion.

Three questions arise from the words of the Gospel, and each of the
interpretations that follow takes one of them up. The first is what the end is
a gathering to. The verse about the body and the eagles is a proverb about
carrion birds, set between the lightning of the coming and the signs in the
sky. St Jerome takes the body, which the Greek calls a fallen thing, for the
Passion of Christ, and the eagles for the saints whose youth is renewed. St
Hilary of Poitiers says that the angels will gather the saints \latin{in loco
passionis}, ``in the place of the Passion.'' St Augustine says that the body is
the one the Lord took up into heaven, and that the spiritual are fed ``as it
were from his body.'' The first interpretation follows them in making the
gathering one to Christ, and finds its beginning in the return the Introit
promises ``from all places'' and in the translation ``into the kingdom'' for
which the Epistle gives thanks.

The second question is to whom the commands of flight are spoken, since
those who hear them on this Sunday are not in Judea. St John Chrysostom
holds that the words concern the Roman war against Jerusalem and are
addressed to the Jews. St Hilary says
that the command to the man on the housetop makes no sense as a rule of
conduct, since nobody can leave a house without coming down from its roof, and
reads it of the person reborn, who must not go down again into the desire of
worldly things. St Jerome adds to the plain senses that every perverse
doctrine standing in the Church is an abomination to be fled. The second
interpretation follows Hilary and Jerome without setting Chrysostom's sense
aside. The worldly desire Hilary forbids is what the Secret asks to be
``delivered from,'' and so the interpretation takes that petition for the
prayer of those who keep his rule.

The third question is what makes a promise certain when no one knows its day.
St John Chrysostom says that ``this generation'' is the generation of
believers, and bids the hearer test the Lord's sayings by what is past and
believe the rest. St Gregory the Great, preaching on Luke's account of the same
discourse, says that the fulfilment of what is past shows the certainty of
what is to come. St Augustine, on the psalm of the Gradual, says that a future
is spoken of as past ``because it is as certain as if it had been done.'' The
third interpretation applies those answers to the past tenses of the chants
and the Epistle and to the Communion's ``believe that you shall receive.''

The English of Scripture quoted in this study is the Douay--Rheims in Bishop
Challoner's revision; the English of the chants and prayers as the Missal
gives them is that of the Lasance \work{New Roman Missal} (1945). The dates
and places of the passages are given in the Scriptural Date and Location
appendix.

\begin{maptable}
Introit &
\latin{Dicit Dominus: Ego cogito cogitationes pacis}; Jer 29:11--12, 14;
verse Ps 84:2 (85:2) &
God's promise to the exiles in Jeremias's letter: ``I think thoughts of
peace, and not of affliction \dots\ I will bring back your captivity from all
places.'' The psalm verse says to God, ``Thou hast turned away the captivity
of Jacob.''\\
Collect &
\latin{Excita, quaesumus, Domine, tuorum fidelium voluntates} &
Asks God to ``arouse \dots\ the wills of Thy faithful,'' that they ``may
the more abundantly partake of Thy mercies.''\\
Epistle &
Col 1:9--14 &
Paul prays that the Colossians be filled with the knowledge of God's will,
walk worthy of him, bear fruit, and be strengthened in patience and
longsuffering with joy, thanking the Father who made them worthy of the lot of
the saints in light, delivered them from the power of darkness and translated
them into the kingdom of his Son.\\
Gradual &
\latin{Liberasti nos, Domine}; Ps 43:8--9 (44:8--9) &
Thanks for deliverance: ``Thou hast delivered us, O Lord, from them that
afflict us.'' The verse promises glory and praise in God's name for ever.\\
Alleluia &
\latin{De profundis clamavi ad te, Domine}; Ps 129:1--2 (130:1--2) &
``From the depths I have cried to Thee, O Lord: Lord, hear my prayer.''\\
Gospel &
Mt 24:15--35 &
The abomination in the holy place and the flight from Judea; the great
tribulation, its days shortened for the elect; false christs; the coming like
lightning; the body and the eagles; the signs in heaven, the Son of man on the
clouds and the gathering of the elect; the parable of the fig tree; this
generation, and the words that do not pass.\\
Offertory &
\latin{De profundis clamavi ad te, Domine}; Ps 129:1--2 (130:1--2) &
The cry from the depths sung again, its first line repeated.\\
Secret &
\latin{Propitius esto, Domine, supplicationibus nostris} &
Asks God to ``accept the offerings and prayers of Thy people'' and to ``turn
all our hearts unto Thee, that, being delivered from earthly desires, we may
pass on to the enjoyments of heaven.''\\
Communion &
\latin{Amen dico vobis, quidquid orantes petitis}; Mk 11:24 &
``Whatsoever you ask when you pray, believe that you shall receive, and it
shall be done unto you.''\\
Postcommunion &
\latin{Concede nobis, quaesumus, Domine} &
Asks that, through the sacrament received, ``whatever is evil in our hearts
may be restored by its gift of healing.''\\
\end{maptable}
